\chapter{Einleitung}
    
Physik ist die Grundlage für das Verständnis technischer Systeme im Maschinenbau und in der Systemtechnik.

Dieses Skript dient als zentrales Medium für den Physikunterricht im 2. Semester für Maschinen- 
und Systemtechniker Stufe HF an der Inovatech. 
Aufbauend auf den Grundkenntnissen der Mathematik vermittelt es die wichtigsten physikalischen Prinzipien, 
die für die Entwicklung und Optimierung moderner Technologien essenziell sind. Von Statik über Dynamik bis hin zur 
Strömungslehre werden die Inhalte möglichst anschaulich und praxisnah aufbereitet.

Wo auch immer möglich werden Beispiele und Aufgaben an reellen physikalischen Systemen gezeigt.
Diese sind meist stark vereinfacht und viele Einflüsse werden vernachlässigt, allerdings ist dies die 
Kunst für alle Stufen der Physik. Die Realität ist beliebig komplex, 
wer sie verständlich und isoliert vereinfachen kann hat einen grossen Vorteil.

Nutzen Sie dieses Skript als Leitfaden, 
um den Einstieg in die Physik zu meistern und Ihre technischen Fähigkeiten zu vertiefen. 
Die behandelten Themen bieten eine solide Basis, die Sie je nach Interesse und Bedarf 
selbstständig vertiefen können, Literatur ist für alle Themengebiete im Überschuss erhältlich. 
Ziel des Unterrichts ist es, Ihnen die wichtigsten Werkzeuge 
an die Hand zu geben, um sich in neue Themengebiete einzuarbeiten und diese erfolgreich anzuwenden.

\begin{figure}[H]    
    \centering          %     left botm right top
    \includegraphics[clip, trim=0cm 0cm 0cm 0cm,width=.3\textwidth]{physics_logo.jpg}
    \label{fig:einleitungPhysicsLogo}
\end{figure}

\newpage

\section{Organisation Unterricht}

Der Unterricht wird über OpenOlat organisiert. Der Unterricht ist als einfache Version des 
Flipped-Classrooms aufgebaut. Auf OpenOlat wird pro Unterrichtsblock ein Themenbereich oder Teile eines
Themenbereichs zugewiesen. Dieses ist als Vorbereitung vor dem jeweiligen Unterrichtsblock vorzubereiten.

In der ersten Lektion des Unterrichtblocks wird in einem kurzen Wrap-up der vorbereitete Stoff wiederholt.
Anschliessend werden allfällige Fragen beantwortet, der Rest des Unterrichts wird für das Lösen von Aufgaben genutzt.

Das Skript kann während des Semesters ändern, bitte achte deshalb darauf, dass du auf
OpenOlat jeweils die aktuellste Version benutzt.

\section{Zu dem Skript}

Das Skript ist in die Themenbereiche aufgeteilt:

\begin{itemize}
    \item Stoffe und Stoffeigenschaften
    \item Statik 1
    \item Kinematik 1
    \item Dynamik
    \item Wellen 1
    \item Bonus
\end{itemize}

Zu jedem Themenbereich sind die Übungen und Lösungen in den Anhängen \ref{app:A_Übungen} und \ref{app:B_Lösungen}.

Das Skript basiert auf vielen verschiedenen Quellen, welche falls möglich referenziert werden.

Das Skript wird von der Lehrperson geschrieben und auf dem aktuellsten Stand gehalten.
Selbstverständlich sind Fehler vorhanden, falls Sie einen bemerken, wird ein Hinweis darauf
dankend angenommen.

\section{Grundlagen}

Dieser Abschnitt behandelt einige Grundlagen, auf welche die Physik aufbaut.

\subsection{Einheiten}

Eine physikalische Größe macht qualitative und quantitative Aussagen über eine messbare Äusserung 
eines physikalischen Zustands oder Vorgangs.
Sie ist das Produkt aus einem Zahlenwert und einer Einheit \cite{boge_physik_2008}.

\begin{infobox}
    physikalische Grösse a = Zahlenwert \{a\} $\cdot$ Einheit [a]
\end{infobox}

Als Grundlage der großen Menge physikalischer
Grössen wurden sieben \textbf{Basisgrössen} festgelegt. Die
Auswahl war willkürlich, mit der einzigen Einschränkung: 
Eine Basisgrösse darf nicht durch
andere Basisgrössen definierbar sein.

Für die Basisgrössen wurden international und gesetzlich 
spezielle Einheiten festgelegt. Sie heissen
\textbf{Basiseinheiten} und sind zugleich Einheiten des
so genannten Internationalen Einheitensystems
(SI-Einheiten). Aus ihnen werden die abgeleiteten
Einheiten gebildet \cite{boge_physik_2008}:

\begin{infobox}
    \begin{minipage}{.3\textwidth}
        \textbf{Basisgrösse} \\
        Länge $l$ \\
        Masse $m$ \\
        Zeit $t$ \\
        Thermod. Temperatur $T$ \\
        Elektrische Stromstärke $I$ \\
        Lichtstärke $I_v$ \\
        Stoffmenge $n$
    \end{minipage}
    \hspace{.1\textwidth}
    \begin{minipage}{.3\textwidth}
        \textbf{Basiseinheit} \\
        Meter ($m$) \\
        Kilogramm ($kg$) \\
        Sekunde ($s$) \\
        Kelvin ($K$) \\
        Ampere ($A$) \\
        Candela ($cd$) \\
        Mol ($mol$)
    \end{minipage}
\end{infobox}

Werden in eine Definitionsgleichung die physikalischen Grössen 
als Produkt von Zahlenwert und
Einheit eingesetzt, entstehen \textbf{abgeleitete Einheiten}. 
Sie sind stets Potenzprodukte von Basiseinheiten 
oder lassen sich auf solche zurückführen \cite{boge_physik_2008}.

\begin{bspbox}
    Die Einheit der Geschwindigkeit ergibt sich aus den Einheiten, durch welche sie 
    berechnet wird. Wenn eine Strecke $s = 1m$ in einer Zeit $t = 1s$ zurückgelegt wird:
    \begin{equation}
        v = \frac{1m}{1s} = 1 \frac{m}{s} = 1 m s^{-1}
    \end{equation}
\end{bspbox}

Die Basiseinheiten m, kg, s, K, A, cd, mol und auch
die abgeleiteten Einheiten mit selbstständigem
Namen (z.B. der Meter $m$) dürfen durch Vorsatzzeichen 
vervielfacht oder unterteilt werden:

\begin{bspbox}
    $1km$ ausgedrückt in $m$, $mm$ und $\mu m$:
    $$
    1km = 1000m = 10^{6}mm = 10^{9}\mu m
    $$
\end{bspbox}

Die gängigsten Vorsatzzeichen:

\begin{minipage}{.3\textwidth}
    $10^{12}$ = Tera ($T$) \\
    $10^9$ = Giga ($G$) \\
    $10^6$ = Mega ($M$) \\
    $10^3$ = Kilo ($k$) \\
    $10^2$ = Hekto ($h$) \\
    $10^1$ = Deka ($da$)
\end{minipage}
\hspace{.1\textwidth}
\begin{minipage}{.3\textwidth}
    $10^{-1}$ = Dezi ($d$) \\
    $10^{-2}$ = Zenti ($c$) \\
    $10^{-3}$ = Milli ($m$) \\
    $10^{-6}$ = Mikro ($\mu$) \\
    $10^{-9}$ = Nano ($n$)  
\end{minipage}

\subsection{Koordinatensysteme}

Das \textbf{kartesische Koordinatensystem} ist das am häufigsten 
verwendete Koordinatensystem. Es besteht aus drei senkrecht aufeinander 
stehenden Achsen: \(x\), \(y\) und \(z\). Jeder Punkt \(P\) im Raum 
wird durch seine Koordinaten \((x, y, z)\) beschrieben.

\begin{center}
    \begin{tikzpicture}[scale=2]
        \draw[->] (0,0,0) -- (2,0,0) node[right]{$x$};
        \draw[->] (0,0,0) -- (0,2,0) node[above]{$y$};
        \draw[->] (0,0,0) -- (0,0,2) node[below left]{$z$};
        \node[circle,fill,inner sep=1pt,label=below:{$P(x,y,z)$}] at (2,2,2) {};
\end{tikzpicture}
\end{center}

Das Zylinderkoordinatensystem oder \textbf{Polarkoordinatensystem} ist 
besonders nützlich für Probleme mit Zylindersymmetrie. 
Ein Punkt \(P\) wird durch die Koordinaten \((\rho, \varphi, z)\) beschrieben, 
wobei \(\rho\) der Abstand zur \(z\)-Achse, \(\varphi\) der Winkel 
in der \(xy\)-Ebene und \(z\) die Höhe über der \(xy\)-Ebene ist.

\begin{center}
\begin{tikzpicture}[scale=2]
    % Axes
    \draw[->] (0,0,0) -- (2,0,0) node[right]{$x$};
    \draw[->] (0,0,0) -- (0,2,0) node[above]{$y$};
    \draw[->] (0,0,0) -- (0,0,2) node[below left]{$z$};

    % Define rho and phi
    \pgfmathsetmacro{\rho}{2}
    \pgfmathsetmacro{\phi}{45}

    % Point P in cylindrical coordinates
    \coordinate (P) at ({\rho*cos(\phi)}, {\rho*sin(\phi)}, 1.5);

    % Draw projection and angle
    \draw[dashed] (0,0,0) -- ({\rho*cos(\phi)}, {\rho*sin(\phi)}, 0) node[midway,below] {$\varrho$};
    \draw (0.5,0,0) arc (0:\phi:0.5) node[midway,right] {$\varphi$};
    \draw[dashed] ({\rho*cos(\phi)}, {\rho*sin(\phi)}, 0) -- (P);

    % Point P
    \node[circle,fill,inner sep=1pt,label=above:{$P(\varrho,\varphi,z)$}] at (P) {};
\end{tikzpicture}
\end{center}